% --- Archon physics formalization source begin ---
% archon:physics
% archon:covers PhyXMiniProblems/problem_phyx_mini_0759.lean
% archon:source-report reports/phyx_mini/problem_phyx_mini_0759.source.json

\chapter{Physics problem phyx\_mini\_0759}
\label{ch:PhyXMiniProblems_problem_phyx_mini_0759}

\paragraph{Problem source.}
\#\# Physical scenario

In figure, a crate of mass \$m = 100\textbackslash{} kg\$ is pushed at constant speed up a frictionless ramp (\$\textbackslash{}theta = 30.0\textasciicircum{}\textbackslash{}circ\$) by a horizontal force \$\textbackslash{}vec\{F\}\$.

\#\# Question

What is the magnitude of the force on the crate from the ramp?

\#\# Answer choices

- A: A: \$0.98 \textbackslash{}times 10\textasciicircum{}3\textbackslash{}\textbackslash{} N\$

- B: B: \$1.02 \textbackslash{}times 10\textasciicircum{}3\textbackslash{}\textbackslash{} N\$

- C: C: \$1.13 \textbackslash{}times 10\textasciicircum{}3\textbackslash{}\textbackslash{} N\$

- D: D: \$1.26 \textbackslash{}times 10\textasciicircum{}3\textbackslash{}\textbackslash{} N\$

\#\# Figure caption (auxiliary; use the image as primary evidence)

The image depicts a wooden crate positioned on an inclined plane. Here's a detailed description:

- **Inclined Plane**: The surface is angled upward to the right, with an angle labeled as \textbackslash{}( \textbackslash{}theta \textbackslash{}) at the bottom left corner.
- **Crate**: A rectangular wooden crate rests on the inclined plane. It’s shown with horizontal slats indicating wood construction. The crate is labeled with the letter \textbackslash{}( m \textbackslash{}), likely representing its mass.
- **Force Vector**: An arrow labeled \textbackslash{}( \textbackslash{}vec\{F\} \textbackslash{}) is shown pointing to the right along the inclined plane, representing the direction and application of an external force.
  
The objects are clearly defined, and the relationships indicate the physics involved with a crate being influenced by a force on an inclined plane.

\#\# Dataset metadata

Dynamics; Physical Model Grounding Reasoning, Spatial Relation Reasoning

\paragraph{Recorded answer/context.}
C: C: \$1.13 \textbackslash{}times 10\textasciicircum{}3\textbackslash{}\textbackslash{} N\$

\paragraph{Figure/image path.}
/root/proposal\_for\_physic/science-mango/phyx\_mini\_run/phyx\_data/test\_image/759.png

\paragraph{Formalization target.}
create a compiling Lean file with sorry bodies at `PhyXMiniProblems/problem\_phyx\_mini\_0759.lean`. The Lean declarations must preserve the physical quantities, dimensions or dimensional roles, figure labels, governing-law hypotheses, and final relation expressed by this problem.
Use Mathlib/Physlib names found through LeanExplore where available. If a physics API is missing, introduce faithful local abstractions rather than scalar placeholder aliases.

\begin{theorem}[Physics formalization target]
\label{thm:physics:phyx_mini_0759:target}
\lean{PhyXMiniProblems.ProblemPhyXMini0759.rampForceMagnitude_is_1960_sqrtThree_over_three_and_matches_C}
\uses{def:physics:phyx-mini-0759:phyxminiproblems-problemphyxmini0759-forceinnewtons, def:physics:phyx-mini-0759:phyxminiproblems-problemphyxmini0759-crateonrampsetup, def:physics:phyx-mini-0759:phyxminiproblems-problemphyxmini0759-matchesproblemdata, def:physics:phyx-mini-0759:phyxminiproblems-problemphyxmini0759-matchessuppliedfigure, def:physics:phyx-mini-0759:phyxminiproblems-problemphyxmini0759-hasphysicalparameters, def:physics:phyx-mini-0759:phyxminiproblems-problemphyxmini0759-satisfiesconstantvelocitykinematics, def:physics:phyx-mini-0759:phyxminiproblems-problemphyxmini0759-satisfiesfrictionlessrampdynamics, def:physics:phyx-mini-0759:phyxminiproblems-problemphyxmini0759-recordedanswerchoice, def:physics:phyx-mini-0759:phyxminiproblems-problemphyxmini0759-isuniqueclosestanswer}
The assigned autoformalize agent should translate this physics problem into Lean declarations in the covered file, with theorem and lemma proof bodies written as `by sorry`.
\end{theorem}
\begin{proof}
This is an autoformalization task, not a proof task. Produce statements that can later be proved by the physics prover without weakening the physical model.
\end{proof}

% --- Archon declaration topology begin ---
\section{Declaration topology}
\begin{definition}[Mass Quantity]
  \label{def:physics:phyx-mini-0759:phyxminiproblems-problemphyxmini0759-massquantity}
  \lean{PhyXMiniProblems.ProblemPhyXMini0759.MassQuantity}
  A nonnegative physical mass
\end{definition}
\begin{proof}
  This is the declared model definition of Mass Quantity.
\end{proof}
\begin{definition}[Force Quantity]
  \label{def:physics:phyx-mini-0759:phyxminiproblems-problemphyxmini0759-forcequantity}
  \lean{PhyXMiniProblems.ProblemPhyXMini0759.ForceQuantity}
  A nonnegative physical force magnitude
\end{definition}
\begin{proof}
  This is the declared model definition of Force Quantity.
\end{proof}
\begin{definition}[Acceleration Quantity]
  \label{def:physics:phyx-mini-0759:phyxminiproblems-problemphyxmini0759-accelerationquantity}
  \lean{PhyXMiniProblems.ProblemPhyXMini0759.AccelerationQuantity}
  A nonnegative physical acceleration magnitude
\end{definition}
\begin{proof}
  This is the declared model definition of Acceleration Quantity.
\end{proof}
\begin{definition}[Signed Acceleration Quantity]
  \label{def:physics:phyx-mini-0759:phyxminiproblems-problemphyxmini0759-signedaccelerationquantity}
  \lean{PhyXMiniProblems.ProblemPhyXMini0759.SignedAccelerationQuantity}
  A signed acceleration component in ramp-adapted coordinates
\end{definition}
\begin{proof}
  This is the declared model definition of Signed Acceleration Quantity.
\end{proof}
\begin{definition}[Speed Quantity]
  \label{def:physics:phyx-mini-0759:phyxminiproblems-problemphyxmini0759-speedquantity}
  \lean{PhyXMiniProblems.ProblemPhyXMini0759.SpeedQuantity}
  A nonnegative physical speed
\end{definition}
\begin{proof}
  This is the declared model definition of Speed Quantity.
\end{proof}
\begin{definition}[mass In Kilograms]
  \label{def:physics:phyx-mini-0759:phyxminiproblems-problemphyxmini0759-massinkilograms}
  \lean{PhyXMiniProblems.ProblemPhyXMini0759.massInKilograms}
  \uses{def:physics:phyx-mini-0759:phyxminiproblems-problemphyxmini0759-massquantity}
  Kilogram readout of a physical mass
\end{definition}
\begin{proof}
  This is the declared model definition of mass In Kilograms.
\end{proof}
\begin{definition}[force In Newtons]
  \label{def:physics:phyx-mini-0759:phyxminiproblems-problemphyxmini0759-forceinnewtons}
  \lean{PhyXMiniProblems.ProblemPhyXMini0759.forceInNewtons}
  \uses{def:physics:phyx-mini-0759:phyxminiproblems-problemphyxmini0759-forcequantity}
  Newton readout of a physical force magnitude
\end{definition}
\begin{proof}
  This is the declared model definition of force In Newtons.
\end{proof}
\begin{definition}[acceleration In Meters Per Second Squared]
  \label{def:physics:phyx-mini-0759:phyxminiproblems-problemphyxmini0759-accelerationinmeterspersecondsquared}
  \lean{PhyXMiniProblems.ProblemPhyXMini0759.accelerationInMetersPerSecondSquared}
  \uses{def:physics:phyx-mini-0759:phyxminiproblems-problemphyxmini0759-accelerationquantity}
  Metres-per-second-squared readout of a nonnegative acceleration
\end{definition}
\begin{proof}
  This is the declared model definition of acceleration In Meters Per Second Squared.
\end{proof}
\begin{definition}[signed Acceleration In Meters Per Second Squared]
  \label{def:physics:phyx-mini-0759:phyxminiproblems-problemphyxmini0759-signedaccelerationinmeterspersecondsquared}
  \lean{PhyXMiniProblems.ProblemPhyXMini0759.signedAccelerationInMetersPerSecondSquared}
  \uses{def:physics:phyx-mini-0759:phyxminiproblems-problemphyxmini0759-signedaccelerationquantity}
  Metres-per-second-squared readout of a signed acceleration component
\end{definition}
\begin{proof}
  This is the declared model definition of signed Acceleration In Meters Per Second Squared.
\end{proof}
\begin{definition}[speed In Meters Per Second]
  \label{def:physics:phyx-mini-0759:phyxminiproblems-problemphyxmini0759-speedinmeterspersecond}
  \lean{PhyXMiniProblems.ProblemPhyXMini0759.speedInMetersPerSecond}
  \uses{def:physics:phyx-mini-0759:phyxminiproblems-problemphyxmini0759-speedquantity}
  Metres-per-second readout of a physical speed
\end{definition}
\begin{proof}
  This is the declared model definition of speed In Meters Per Second.
\end{proof}
\begin{definition}[degrees To Radians]
  \label{def:physics:phyx-mini-0759:phyxminiproblems-problemphyxmini0759-degreestoradians}
  \lean{PhyXMiniProblems.ProblemPhyXMini0759.degreesToRadians}
  Convert an angle stated in degrees to its dimensionless radian value
\end{definition}
\begin{proof}
  This is the declared model definition of degrees To Radians.
\end{proof}
\begin{definition}[Contact Model]
  \label{def:physics:phyx-mini-0759:phyxminiproblems-problemphyxmini0759-contactmodel}
  \lean{PhyXMiniProblems.ProblemPhyXMini0759.ContactModel}
  Idealized contact between the crate and the ramp
\end{definition}
\begin{proof}
  This is the declared model definition of Contact Model.
\end{proof}
\begin{definition}[Ramp Motion Direction]
  \label{def:physics:phyx-mini-0759:phyxminiproblems-problemphyxmini0759-rampmotiondirection}
  \lean{PhyXMiniProblems.ProblemPhyXMini0759.RampMotionDirection}
  Direction of the crate's translational motion along the ramp
\end{definition}
\begin{proof}
  This is the declared model definition of Ramp Motion Direction.
\end{proof}
\begin{definition}[Speed Regime]
  \label{def:physics:phyx-mini-0759:phyxminiproblems-problemphyxmini0759-speedregime}
  \lean{PhyXMiniProblems.ProblemPhyXMini0759.SpeedRegime}
  Whether the scalar speed changes during the depicted motion
\end{definition}
\begin{proof}
  This is the declared model definition of Speed Regime.
\end{proof}
\begin{definition}[Force Direction]
  \label{def:physics:phyx-mini-0759:phyxminiproblems-problemphyxmini0759-forcedirection}
  \lean{PhyXMiniProblems.ProblemPhyXMini0759.ForceDirection}
  Distinguished force directions in the vertical cross-section
\end{definition}
\begin{proof}
  This is the declared model definition of Force Direction.
\end{proof}
\begin{definition}[Ramp Orientation]
  \label{def:physics:phyx-mini-0759:phyxminiproblems-problemphyxmini0759-ramporientation}
  \lean{PhyXMiniProblems.ProblemPhyXMini0759.RampOrientation}
  Qualitative orientation of the ramp in the supplied image
\end{definition}
\begin{proof}
  This is the declared model definition of Ramp Orientation.
\end{proof}
\begin{definition}[Figure Component]
  \label{def:physics:phyx-mini-0759:phyxminiproblems-problemphyxmini0759-figurecomponent}
  \lean{PhyXMiniProblems.ProblemPhyXMini0759.FigureComponent}
  Physical elements explicitly visible in the supplied image
\end{definition}
\begin{proof}
  This is the declared model definition of Figure Component.
\end{proof}
\begin{definition}[Figure Label]
  \label{def:physics:phyx-mini-0759:phyxminiproblems-problemphyxmini0759-figurelabel}
  \lean{PhyXMiniProblems.ProblemPhyXMini0759.FigureLabel}
  Symbols printed in the supplied image
\end{definition}
\begin{proof}
  This is the declared model definition of Figure Label.
\end{proof}
\begin{definition}[Supplied Figure]
  \label{def:physics:phyx-mini-0759:phyxminiproblems-problemphyxmini0759-suppliedfigure}
  \lean{PhyXMiniProblems.ProblemPhyXMini0759.SuppliedFigure}
  \uses{def:physics:phyx-mini-0759:phyxminiproblems-problemphyxmini0759-forcedirection, def:physics:phyx-mini-0759:phyxminiproblems-problemphyxmini0759-ramporientation, def:physics:phyx-mini-0759:phyxminiproblems-problemphyxmini0759-figurecomponent, def:physics:phyx-mini-0759:phyxminiproblems-problemphyxmini0759-figurelabel}
  A structured readout of the primary image
\end{definition}
\begin{proof}
  This is the declared model definition of Supplied Figure.
\end{proof}
\begin{definition}[Crate On Ramp Setup]
  \label{def:physics:phyx-mini-0759:phyxminiproblems-problemphyxmini0759-crateonrampsetup}
  \lean{PhyXMiniProblems.ProblemPhyXMini0759.CrateOnRampSetup}
  \uses{def:physics:phyx-mini-0759:phyxminiproblems-problemphyxmini0759-massquantity, def:physics:phyx-mini-0759:phyxminiproblems-problemphyxmini0759-forcequantity, def:physics:phyx-mini-0759:phyxminiproblems-problemphyxmini0759-accelerationquantity, def:physics:phyx-mini-0759:phyxminiproblems-problemphyxmini0759-signedaccelerationquantity, def:physics:phyx-mini-0759:phyxminiproblems-problemphyxmini0759-speedquantity, def:physics:phyx-mini-0759:phyxminiproblems-problemphyxmini0759-contactmodel, def:physics:phyx-mini-0759:phyxminiproblems-problemphyxmini0759-rampmotiondirection, def:physics:phyx-mini-0759:phyxminiproblems-problemphyxmini0759-speedregime, def:physics:phyx-mini-0759:phyxminiproblems-problemphyxmini0759-forcedirection, def:physics:phyx-mini-0759:phyxminiproblems-problemphyxmini0759-suppliedfigure}
  The crate-ramp apparatus and its unknown force magnitudes and acceleration components. In particular, rampForceMagnitude is an independent physical quantity; it is not defined from a displayed answer or from the desired closed form
\end{definition}
\begin{proof}
  This is the declared model definition of Crate On Ramp Setup.
\end{proof}
\begin{definition}[Matches Problem Data]
  \label{def:physics:phyx-mini-0759:phyxminiproblems-problemphyxmini0759-matchesproblemdata}
  \lean{PhyXMiniProblems.ProblemPhyXMini0759.MatchesProblemData}
  \uses{def:physics:phyx-mini-0759:phyxminiproblems-problemphyxmini0759-massinkilograms, def:physics:phyx-mini-0759:phyxminiproblems-problemphyxmini0759-accelerationinmeterspersecondsquared, def:physics:phyx-mini-0759:phyxminiproblems-problemphyxmini0759-crateonrampsetup}
  Problem-statement data and the standard gravitational calibration used by the displayed numerical choices. No value of the ramp force occurs here
\end{definition}
\begin{proof}
  This is the declared model definition of Matches Problem Data.
\end{proof}
\begin{definition}[Matches Supplied Figure]
  \label{def:physics:phyx-mini-0759:phyxminiproblems-problemphyxmini0759-matchessuppliedfigure}
  \lean{PhyXMiniProblems.ProblemPhyXMini0759.MatchesSuppliedFigure}
  \uses{def:physics:phyx-mini-0759:phyxminiproblems-problemphyxmini0759-crateonrampsetup}
  Primary-image readout: the ramp rises to the right, the crate is labelled m, θ labels the incline, and the arrow labelled F is horizontal
\end{definition}
\begin{proof}
  This is the declared model definition of Matches Supplied Figure.
\end{proof}
\begin{definition}[Has Physical Parameters]
  \label{def:physics:phyx-mini-0759:phyxminiproblems-problemphyxmini0759-hasphysicalparameters}
  \lean{PhyXMiniProblems.ProblemPhyXMini0759.HasPhysicalParameters}
  \uses{def:physics:phyx-mini-0759:phyxminiproblems-problemphyxmini0759-massinkilograms, def:physics:phyx-mini-0759:phyxminiproblems-problemphyxmini0759-forceinnewtons, def:physics:phyx-mini-0759:phyxminiproblems-problemphyxmini0759-accelerationinmeterspersecondsquared, def:physics:phyx-mini-0759:phyxminiproblems-problemphyxmini0759-speedinmeterspersecond, def:physics:phyx-mini-0759:phyxminiproblems-problemphyxmini0759-degreestoradians, def:physics:phyx-mini-0759:phyxminiproblems-problemphyxmini0759-crateonrampsetup}
  Positivity and acute-angle conditions selecting the physical branch
\end{definition}
\begin{proof}
  This is the declared model definition of Has Physical Parameters.
\end{proof}
\begin{definition}[Satisfies Constant Velocity Kinematics]
  \label{def:physics:phyx-mini-0759:phyxminiproblems-problemphyxmini0759-satisfiesconstantvelocitykinematics}
  \lean{PhyXMiniProblems.ProblemPhyXMini0759.SatisfiesConstantVelocityKinematics}
  \uses{def:physics:phyx-mini-0759:phyxminiproblems-problemphyxmini0759-signedaccelerationinmeterspersecondsquared, def:physics:phyx-mini-0759:phyxminiproblems-problemphyxmini0759-crateonrampsetup}
  For straight-line motion up the ramp at constant speed, both ramp-adapted components of the crate's acceleration vanish
\end{definition}
\begin{proof}
  This is the declared model definition of Satisfies Constant Velocity Kinematics.
\end{proof}
\begin{definition}[weight In Newtons]
  \label{def:physics:phyx-mini-0759:phyxminiproblems-problemphyxmini0759-weightinnewtons}
  \lean{PhyXMiniProblems.ProblemPhyXMini0759.weightInNewtons}
  \uses{def:physics:phyx-mini-0759:phyxminiproblems-problemphyxmini0759-massinkilograms, def:physics:phyx-mini-0759:phyxminiproblems-problemphyxmini0759-accelerationinmeterspersecondsquared, def:physics:phyx-mini-0759:phyxminiproblems-problemphyxmini0759-crateonrampsetup}
  Gravitational-force magnitude m g, read in newtons
\end{definition}
\begin{proof}
  This is the declared model definition of weight In Newtons.
\end{proof}
\begin{definition}[Satisfies Frictionless Ramp Dynamics]
  \label{def:physics:phyx-mini-0759:phyxminiproblems-problemphyxmini0759-satisfiesfrictionlessrampdynamics}
  \lean{PhyXMiniProblems.ProblemPhyXMini0759.SatisfiesFrictionlessRampDynamics}
  \uses{def:physics:phyx-mini-0759:phyxminiproblems-problemphyxmini0759-massinkilograms, def:physics:phyx-mini-0759:phyxminiproblems-problemphyxmini0759-forceinnewtons, def:physics:phyx-mini-0759:phyxminiproblems-problemphyxmini0759-signedaccelerationinmeterspersecondsquared, def:physics:phyx-mini-0759:phyxminiproblems-problemphyxmini0759-degreestoradians, def:physics:phyx-mini-0759:phyxminiproblems-problemphyxmini0759-crateonrampsetup, def:physics:phyx-mini-0759:phyxminiproblems-problemphyxmini0759-weightinnewtons}
  Newton's second law resolved tangent to the ramp and along the outward ramp normal. A horizontal push has tangential component F cos θ and inward-normal component F sin θ; gravity contributes inward components mg sin θ and mg cos θ. The frictionless contact force is normal to the ramp. These are governing component balances. They do not state the requested normal-force value or any displayed answer
\end{definition}
\begin{proof}
  This is the declared model definition of Satisfies Frictionless Ramp Dynamics.
\end{proof}
\begin{lemma}[ramp Force eq weight div cosine]
  \label{lem:physics:phyx-mini-0759:phyxminiproblems-problemphyxmini0759-rampforce-eq-weight-div-cosine}
  \lean{PhyXMiniProblems.ProblemPhyXMini0759.rampForce_eq_weight_div_cosine}
  \uses{def:physics:phyx-mini-0759:phyxminiproblems-problemphyxmini0759-forceinnewtons, def:physics:phyx-mini-0759:phyxminiproblems-problemphyxmini0759-degreestoradians, def:physics:phyx-mini-0759:phyxminiproblems-problemphyxmini0759-crateonrampsetup, def:physics:phyx-mini-0759:phyxminiproblems-problemphyxmini0759-matchesproblemdata, def:physics:phyx-mini-0759:phyxminiproblems-problemphyxmini0759-hasphysicalparameters, def:physics:phyx-mini-0759:phyxminiproblems-problemphyxmini0759-satisfiesconstantvelocitykinematics, def:physics:phyx-mini-0759:phyxminiproblems-problemphyxmini0759-weightinnewtons, def:physics:phyx-mini-0759:phyxminiproblems-problemphyxmini0759-satisfiesfrictionlessrampdynamics}
  Eliminating the unknown horizontal push from the two zero-acceleration component balances gives N = mg / cos θ
\end{lemma}
\begin{proof}
  The conclusion follows from the governing relations and figure readouts named in the statement of ramp Force eq weight div cosine.
\end{proof}
\begin{definition}[Answer Choice]
  \label{def:physics:phyx-mini-0759:phyxminiproblems-problemphyxmini0759-answerchoice}
  \lean{PhyXMiniProblems.ProblemPhyXMini0759.AnswerChoice}
  Labels of the four displayed answers
\end{definition}
\begin{proof}
  This is the declared model definition of Answer Choice.
\end{proof}
\begin{definition}[force In Newtons]
  \label{def:physics:phyx-mini-0759:phyxminiproblems-problemphyxmini0759-answerchoice-forceinnewtons}
  \lean{PhyXMiniProblems.ProblemPhyXMini0759.AnswerChoice.forceInNewtons}
  \uses{def:physics:phyx-mini-0759:phyxminiproblems-problemphyxmini0759-forceinnewtons, def:physics:phyx-mini-0759:phyxminiproblems-problemphyxmini0759-answerchoice}
  Force magnitudes in newtons printed beside the answer labels
\end{definition}
\begin{proof}
  This is the declared model definition of force In Newtons.
\end{proof}
\begin{definition}[recorded Answer Choice]
  \label{def:physics:phyx-mini-0759:phyxminiproblems-problemphyxmini0759-recordedanswerchoice}
  \lean{PhyXMiniProblems.ProblemPhyXMini0759.recordedAnswerChoice}
  \uses{def:physics:phyx-mini-0759:phyxminiproblems-problemphyxmini0759-answerchoice}
  The answer label recorded in the supplied dataset metadata
\end{definition}
\begin{proof}
  This is the declared model definition of recorded Answer Choice.
\end{proof}
\begin{definition}[Is Unique Closest Answer]
  \label{def:physics:phyx-mini-0759:phyxminiproblems-problemphyxmini0759-isuniqueclosestanswer}
  \lean{PhyXMiniProblems.ProblemPhyXMini0759.IsUniqueClosestAnswer}
  \uses{def:physics:phyx-mini-0759:phyxminiproblems-problemphyxmini0759-forceinnewtons, def:physics:phyx-mini-0759:phyxminiproblems-problemphyxmini0759-crateonrampsetup, def:physics:phyx-mini-0759:phyxminiproblems-problemphyxmini0759-answerchoice}
  A displayed choice is uniquely closest to the derived ramp-force value
\end{definition}
\begin{proof}
  This is the declared model definition of Is Unique Closest Answer.
\end{proof}
% --- Archon declaration topology end ---
% --- Archon physics formalization source end ---
